\cotefile{157-3}
\scanpage{297}
\mspage{247}

\uline{Corollaire}. Description \unsure{du|d'un} foncteur
\[
(\xi_1,\xi_2)\mathrel{\unsure{\vdash|\mapsto}}\xi_1\times_S\xi_2
\colon \mathrm{Ho}^{\unsure{\mathrm{rel}|\mathrm{le}}}_W(S)\times\mathrm{Ho}^{\mathrm{rel}}_W(S)
\mathrel{\unsure{-|\to}}\mathrm{Ho}^{\mathrm{rel}}_W(S)
\]

Ce \unsure{foncteur} est bel et bien un produit cartésien dans $\mathrm{Ho}^{\mathrm{rel}}_W(S)$\inserted{mais c'est un cas particulier du cas de $\mathrm{Ho}_W$ (\unsure{c.-à-d.|i.e.|ci dt.} des produits binaires), en remplaçant $M$ par $M_{/S}$}, \unsure{i.e.} les produits finis existent.$^{*}$

% slanted
\margin{$^{*}$ Dire aussi que le foncteur $f^*\colon\mathrm{Ho}^{\mathrm{rel}}_W(S)\longrightarrow\mathrm{Ho}^{\mathrm{rel}}_W(S')$ (pour $f\colon S'\longrightarrow S$) \unsure{commute} aux produits finis.}

Pour prouver ce dernier point, \unsure{posant|posons}
\[
\mathscr H=\mathrm{Ho}^{\mathrm{rel}}_W(S)
\]
il faut montrer que les foncteurs
\diagram{\begin{tikzcd}[ampersand replacement=\&]
\mathscr H \arrow[r,shift left,"\delta"] \& \mathscr H\times\mathscr H \arrow[l,shift left,"\pi"]
\end{tikzcd}}
sont adjoints, \unsure{où|avec}
\[
\begin{aligned}
\delta &= \operatorname{diag}_{\mathscr H}\colon &\xi&\longmapsto(\xi,\xi)\\
\pi &= \text{``produit''}\colon &(\xi_1,\xi_2)&\mathrel{\unsure{\vdash|\mapsto}}\xi_1\times_S\xi_2.
\end{aligned}
\]

On va définir les hom\unsure{.|} d'adjonction
\[
\begin{aligned}
\mathrm{id}_{\mathscr H}\mathrel{\overset{\alpha}{\unsure{\longleftrightarrow|\longrightarrow}}}\pi\delta
&\quad\text{i.e.}\quad \xi\xrightarrow{\alpha_\xi}\xi\times_S\xi\\
\mathrm{id}_{\mathscr H\times\mathscr H}\xleftarrow{\beta}\delta\pi
&\quad\text{i.e.}\quad(\xi_1,\xi_2)\xleftarrow{\beta_{\unsure{\xi_1|1},\xi_2}}(\xi_1\times_S\xi_2,\xi_1\times_S\xi_2)
\end{aligned}
\]

Pour ceci, on décrit $\mathscr H$ comme
\[
\mathscr H=(\underline{\mathrm{Fib}}_W S)(W_S)^{-1}
\]
et on note que $\delta,\pi$ proviennent de foncteurs adjoints au niveau de $\mathscr H_0=\underline{\mathrm{Fib}}_W S$
\diagram{\begin{tikzcd}[ampersand replacement=\&]
\mathscr H_0 \arrow[r,shift left,"\delta_0"] \& \mathscr H_0\times\mathscr H_0 \arrow[l,shift left,"\pi_0"]
\end{tikzcd}}
\deleted{\unsure{Ce|De}}\unsure{les}quels donnent lieu : $\alpha_0,\beta_0$
\[
\begin{aligned}
\mathrm{id}_{\mathscr H_0}&\xleftarrow{\alpha_0}\pi_0\delta_0\\
\mathrm{id}_{\mathscr H_0\times\mathscr H_0}&\mathrel{\overset{\beta_0}{\unsure{\longrightarrow|-}}}\delta_0\pi_0 .
\end{aligned}
\]
% archive sheet number: 297
\scanpage{298}
\mspage{248}

On définit $\alpha,\beta$ par \unsure{passage} aux cat. de fractions pour $\alpha_0,\beta_0$. Les compatibilités, étant vérifiées pour $\alpha_0,\beta_0$, le sont aussi pour $\alpha,\beta$. \unsure{Qed.|Qed}

Je \deleted{\unsure{dé}}dégage ici le

\uline{Lemme} \unsure{1|8} : Soit
\diagram{\begin{tikzcd}[ampersand replacement=\&]
\mathcal X\arrow[r,shift left,"f"] \&\mathcal Y\arrow[l,shift left,"g"]
\end{tikzcd}}
un couple de foncteurs adjoints, et soient $\Sigma,\Sigma'$ des \unsure{localiseurs|localisants} dans $\mathcal X,\mathcal Y$ tels que $f(\Sigma)\subset\Sigma'$, $g(\Sigma')\subset\Sigma$, de sorte qu'on a un couple de foncteurs
\diagram{\begin{tikzcd}[ampersand replacement=\&]
\deleted{\unsure{\mathcal X}}\ \mathcal X\Sigma^{-1}\arrow[r,shift left,dash,"\bar f"]\&\mathcal Y{\Sigma'}^{-1}\arrow[l,shift left,"\bar g"]
\end{tikzcd}.}

Alors $\bar f,\bar g$ sont adjoints, les flèches d'adjonction étant \unsure{déduites} de celles pour $f,g$.

\uline{Corollaire} Soit $(M,W)$ une cat. de modèles, \deleted{\unsure{avec}} supposons que dans $M$ les produits binaires (resp. les produits finis) existent, et que \deleted{le produit \unsure{cartésien} de} pour $u_1\colon X_1\mathrel{\unsure{-|\to}} X'_1$ et $u_2\colon X_2\to X'_2$, si $u_1,u_2\in W$, alors $u_1\times u_2\colon X_1\times X_2\to X'_1\times X'_2$ est dans $W$. Alors dans $MW^{-1}$ les produits binaires (resp. produits finis) existent, et le foncteur \inserted{localisation} $M\mathrel{\unsure{-|\to}} MW^{-1}$ y commute.

On applique le lemme à
\diagram{\begin{tikzcd}[ampersand replacement=\&]
M\arrow[r,shift left,"{\deleted{\unsure{\mathrm{diag}}}\delta}"]\& M\times M\arrow[l,shift left,"\pi"]
\end{tikzcd}}
et $W,W\times W$, et pour l'objet final,
\diagram{\begin{tikzcd}[ampersand replacement=\&]
M\arrow[r,shift left,"{\unsure{i}}"]\& e\arrow[l,shift left,"\bullet"]
\end{tikzcd}}
et $W,\deleted{\unsure{\{\mathrm{id}_e\}}}W_0=\{\mathrm{id}_e\}$
% archive sheet number: 298
\scanpage{299}
\mspage{249}

Je vais maintenant construire le foncteur
\[
i_*\colon\mathrm{Ho}_W(\Psi)\longrightarrow\mathrm{Ho}_W(\Psi^{\sim})
\]

Il me faut \deleted{\unsure{donner}} d'abord prouver le

\uline{Lemme 1} Soit $\mathscr F$ la sous-catégorie pleine de $M(\Psi)$, catégorie des diagrammes
\diagram{\begin{tikzcd}[ampersand replacement=\&,row sep=small,column sep=small]
\& X_1\arrow[dl,"p_1"']\\
S\\
\& X_2\arrow[ul,"p_2"]
\end{tikzcd}}
dans $S$, formée de ceux pour lesquels $p_1,p_2$ sont dans $\mathrm{Fib}_W$. Soit $W_{\mathscr F}$ le localiseur induit sur $\mathscr F$ par $W(\Psi)$\unsure{,} ($W$-équivalences terme à terme). Alors
\[
\mathscr F W_{\mathscr F}^{-1}\xrightarrow{\approx}M(\Psi)W(\Psi)^{-1}
\bigl(\mathrel{\overset{\mathrm{def}}{=}}\mathrm{Ho}_W(\Psi)\bigr)
\]

Nous allons prouver ce lemme plus bas.

En l'admettant, pour construire $i_*$, il suffit de construire
\[
\mathscr F\longrightarrow\mathrm{Ho}_W(\Psi^{\sim})
\]
qui \deleted{\unsure{tr}}applique $W_{\mathscr F}$ \deleted{dans} en des iso. Je \unsure{prendrai|prendra}, bien sûr
\diagram{\begin{tikzcd}[ampersand replacement=\&,row sep=small,column sep=small]
\& X_1\arrow[dl]\& \& \& X_1\arrow[dl,dash] \&\\
S\& \& {\longmapsto}\&S\& \&X_1\times_S X_2\arrow[ul]\arrow[dl]\\
\& X_2\arrow[ul]\&\&\& X_2\arrow[ul]\&
\end{tikzcd}}
\[
\mathscr F\longrightarrow\mathrm{Ho}_W(\Psi)
\]
Vérifions qu'il transforme \deleted{\unsure{$W_{\mathscr F}$|fl}}\unsure{flèches|flèche} de $W_{\mathscr F}$ en iso.
% archive sheet number: 299
\scanpage{300}
\mspage{250}
\diagram{\begin{tikzcd}[ampersand replacement=\&,row sep=small,column sep=small]
\& X_1\times_S X_2\arrow[dl]\arrow[ddr,dash]\&\&\&X'_1\times_{S'}X'_2\arrow[lll,dash,"{f_3}"']\arrow[dl,dash]\arrow[ddr]\&\\
X_1\&\&\&X'_1\arrow[lll,"{f_1}"',near start]\arrow[ddr,dash]\&\&\\
\&\& X_2\arrow[dl]\&\&\&X'_2\arrow[lll,dash,"{f_2}"]\arrow[dl]\\
\&S\arrow[uul]\&\&\&S'\arrow[lll,"{f_0}"]\&
\end{tikzcd}}

Il faut prouver que si $f_0,f_1,f_2\in W$, alors $f_3\in W$. Mais considérons l'image inverse du carré cartésien $Q$ par $f_0\colon S'\mathrel{\unsure{-|\to}}S$, d'où un carré cartésien $Q_{S'}$ et une factorisation
\diagram{\begin{tikzcd}[ampersand replacement=\&]
Q\& Q_{S'}\arrow[l,"{\widetilde f}"']\& Q'\arrow[l,"g"']\arrow[ll,bend left=18,"f"]
\end{tikzcd}}
de $Q'\mathrel{-}Q$. Comme $f_0\in W$, et que $X_1,X_2,X_1\times_S X_2$ sont dans $\underline{\mathrm{Fib}}_W S$, les flèches horizontales \unsure{réalisant} $Q_{S'}\mathrel{-}Q$ sont toutes dans $W$. Ceci, plus le fait que $f_1,\unsure{f|f_2}\in W$, \unsure{impliquent} que dans le morphisme $g$, les flèches $g_1,g_2$ sont dans $W$. Il en résulte que \unsure{la|le} flèche
\[
g_3\colon X'_1\times_{S'}X'_2\mathrel{\unsure{-|\to}}(X_1\times_S X_2)\times_S S'
\]
l'est aussi (théorie du produit dans $\underline{\mathrm{Fib}}_W S'$ et $\mathrm{Ho}^{\unsure{\mathrm{rel}|\mathrm{le}}}_W(S')$), donc $f_3=\widetilde f_3g_3\in W$, \unsure{qed}.
% archive sheet number: 300
\scanpage{301}
\mspage{251}

Notons qu'on a un iso. canonique
\[
\alpha\colon i^*i_*\xrightarrow{\sim}\mathrm{id}_{\mathrm{Ho}_W(\unsure{\Phi|\Psi})}\tag*{\makebox[0pt][r]{\makebox[\displaywidth][l]{$(\unsure{*|x})$}}}
\]

J'ai envie de prouver que les foncteurs
\diagram{\begin{tikzcd}[ampersand replacement=\&]
\mathrm{Ho}_W(\Phi)\arrow[r,shift right,"{i_*}"']\&\mathrm{Ho}_W(\Phi^{\sim})\arrow[l,shift right,"{i^*}"']
\end{tikzcd}}
sont adjoints, et que $\unsure{\alpha}$ est une des flèches d'adjonction. Il faut définir l'autre flèche d'adjonction
\[
\beta\colon i_*i^*\longleftarrow\mathrm{id}_{\mathrm{Ho}_W(\Phi^{\sim})}.\tag*{\makebox[0pt][r]{\makebox[\displaywidth][l]{$(**)$}}}
\]

Mais pour ceci, je dois utiliser le

\uline{Lemme 2} Soit $\mathscr G$ la sous-catégorie pleine de $M(\Phi^{\sim})$, formée des
\diagram{\begin{tikzcd}[ampersand replacement=\&]
X_1\arrow[d,"{p_1}"']\& X_3\arrow[l,"{q_2}"']\arrow[d,"{q_{\deleted{\unsure{2}}1}}"]\\
S\&X_2\arrow[l,"{p_2}"]
\end{tikzcd}}
tels que $p_1,p_2\in\mathrm{Fib}_W$, et soit $W_{\mathscr G}$ le localiseur induit par $W(\Phi^{\sim})$ ($W$-équivalences terme à terme). Alors l'inclusion $\mathscr G\mathrel{\unsure{-|\to}}M(\Phi^{\sim})$ induit une équivalence
\[
\mathscr G W_{\mathscr G}^{-1}\xrightarrow{\approx}M(\Phi^{\sim})W(\Phi^{\sim})^{-1}
\mathrel{\overset{\mathrm{def}}{=}}\mathrm{Ho}_W(\Phi^{\sim})
\]

Je vais pour le moment admettre ce
% archive sheet number: 301
\scanpage{302}
\mspage{252}
lemme, \unsure{tout} comme le lemme 1\inserted{(p.~249)}. Je vais considérer $i^*,i_*$ comme des flèches entre les catégories $\mathscr F W_{\mathscr F}^{-1},\mathscr G W_{\mathscr G}^{-1}$. Elles sont induites par les flèches
\diagram{\begin{tikzcd}[ampersand replacement=\&]
\deleted{\unsure{\mathscr G}}\ \mathscr F\arrow[r,shift right,"{i_{0*}}"']\&\mathscr G\arrow[l,shift right,"{i_0^*}"']
\end{tikzcd}}
entre $\mathscr F$ et $\mathscr G$ eux-mêmes
\diagram{\begin{tikzcd}[ampersand replacement=\&,column sep=small,row sep=small]
{i_0^*\colon\mathscr G\mathrel{\unsure{-|\to}}\mathscr F\colon}\&X_1\arrow[d]\&X_3\arrow[l]\arrow[d]\&{\longmapsto}\&X_1\arrow[d]\&\\
\&S\&X_2\arrow[l]\&\&S\&X_2\arrow[l]
\end{tikzcd}}
\diagram{\begin{tikzcd}[ampersand replacement=\&,column sep=small,row sep=small]
{i_{0*}\colon\mathscr F\to\mathscr G\colon}\&X_1\arrow[d]\&\&{\longmapsto}\&X_1\arrow[d]\&X_1\times_S X_2\arrow[l,arrows={-{Bar}}]\arrow[d]\\
\&S\&X_2\arrow[l]\&\&S\&X_2\arrow[l]
\end{tikzcd}.}

Ces deux foncteurs sont adjoints l'un de l'autre, et compatibles aux localiseurs $W_{\mathscr F},W_{\mathscr G}$. Donc \deleted{\unsure{il}}les foncteurs \deleted{\unsure{obtenus}} déduits de $i_0^*,i_{0*}$ par passage aux catégories de fractions sont également adjoints (cf lemme p.~248).

Ainsi, modulo la dém. des lemmes 1 et \unsure{2}, on a prouvé ceci :
% archive sheet number: 302
\scanpage{303}
\mspage{253}

\uline{Théorème} Soit $(M,W)$ une catégorie de modèles $(W\subset\mathrm{Fl}\,M)$ satisfaisant l'axiome de factorisation \unsure{gauche} (p.~240), donnant \inserted{ainsi} lieu à une bonne théorie des \deleted{\unsure{la}} fibres $W$-homotopiques. \unsure{Soient|Soit}
\diagram{\begin{tikzcd}[ampersand replacement=\&,row sep=small,column sep=small]
\&\&1\&\&\&\&1\arrow[dr,leftrightarrow]\&\\
{\Psi=}\&0\arrow[ur,leftrightarrow]\arrow[dr,leftrightarrow]\&\&{}\arrow[r,"i"]\&{\Psi^{\sim}=}\&0\arrow[ur]\arrow[dr,leftrightarrow]\&\&3\\
\&\&2\&\&\&\&2\arrow[ur]\&
\end{tikzcd}}
(immersion ouverte dans $\mathrm{Cat}$), et considérons le foncteur restriction
\[
i^*\colon\mathrm{Ho}_W(\Psi^{\sim})\longrightarrow\mathrm{Ho}_W(\Psi).
\]

Ce foncteur admet un adjoint à droite
\[
\unsure{j|i}_*\colon\mathrm{Ho}_W(\Psi)\longrightarrow\mathrm{Ho}_W(\Psi^{\sim}).
\]

Ce foncteur $i_*$ se ``calcule'' ainsi : sa valeur en
\diagram{\begin{tikzcd}[ampersand replacement=\&,row sep=small,column sep=small]
X_1\arrow[d,"{p_1}"']\&\deleted{\unsure{X_3}}\arrow[l,draw opacity=.45,"{\deleted{\leftarrow}}" description]\arrow[d,draw opacity=.45,"{\deleted{\downarrow}}" description]\\
S\& X_2\arrow[l,"{p_2}"]
\end{tikzcd}}
est canoniquement \unsure{isom|iso} (dans $\mathrm{Ho}_W(\Psi^{\sim})$) \unsure{:|à}
\diagram{\begin{tikzcd}[ampersand replacement=\&,row sep=small,column sep=small]
X_1\arrow[d]\&X_1\times_S X_2\arrow[l]\arrow[d]\\
S\&X_2\arrow[l]
\end{tikzcd},}
si on suppose que $p_1$ ou $p_2$ est une $W$-fibration. Dans le cas général elle se calcule en factorisant $p_1$ \deleted{\unsure{en $X$}} en
\[
X_1\mathrel{\overset{i_1}{\unsure{-|\to}}}\overline X_1\xrightarrow{p'_1}S,\qquad i_1\in W,\ \unsure{f|p}'_1\in\mathrm{Fib}_W,
\]
et en prenant
\diagram{\begin{tikzcd}[ampersand replacement=\&,row sep=small,column sep=small]
\overline X_1\arrow[d]\&\overline X_1\times_S X_2\arrow[l]\arrow[d]\\
S\&X_2\arrow[l]
\end{tikzcd},}
ou en
% archive sheet number: 303
\scanpage{304}
\mspage{254}
procédant de façon symétrique, ou en factorisant simultanément $p_1$ et $p_2$, et prenant
\diagram{\begin{tikzcd}[ampersand replacement=\&]
\overline X_1\arrow[d]\&\overline X_1\times_S\overline X_2\arrow[l,arrows={-{Bar}}]\arrow[d,dash]\\
S\&\overline X_2\arrow[l]
\end{tikzcd}.}

La loi fonctorielle s'obtient à l'aide du lemme 1 p.~249, qui permet de se borner aux arguments
\diagram{\begin{tikzcd}[ampersand replacement=\&,row sep=small,column sep=small]
X_1\arrow[d,"{p_1}"']\&\\S\&X_2\arrow[l,arrows={-{Bar}},"{p_2}"]
\end{tikzcd}}
tels que $p_1,p_2\in\mathrm{Fib}_W$, de sorte que l'on \unsure{donne|trouve} un foncteur
\diagram{\begin{tikzcd}[ampersand replacement=\&,row sep=small,column sep=small]
X_1\arrow[d]\&\&{\longmapsto}\&X_1\arrow[d]\&X_1\times_S X_2\arrow[l]\arrow[d]\\
S\&X_2\arrow[l]\&\&S\&X_2\arrow[l]
\end{tikzcd}}
compatible aux localiseurs, et passant aux catégories de fractions $\mathrm{Ho}_W(\Psi)$ et $\mathrm{Ho}_W(\Psi^{\sim})$ respectivement.

\uline{NB} En procédant de façon similaire, plus simple (sans besoin d'un recours au lemme 1), on trouve l'existence d'un \uline{adjoint à gauche} $\unsure{i|j}_!$ de $i^*\colon\mathrm{Ho}_W(\Psi^{\sim})\to\mathrm{Ho}_W(\Psi)$, sous réserve de l'existence d'un objet initial $\unsure{\varnothing|\mathscr O|\varphi}$ dans $M$. \unsure{Cet|Ce} foncteur $\unsure{j|i}_!$ est donné par
% archive sheet number: 304
\scanpage{305}
\mspage{255}
\diagram{\(\displaystyle
i_!\left(\begin{tikzcd}[ampersand replacement=\&,row sep=small,column sep=small]
X_1\arrow[d]\&\\S\&X_2\arrow[l]
\end{tikzcd}\right)=
\left(\begin{tikzcd}[ampersand replacement=\&,row sep=small,column sep=small]
X_1\arrow[d]\&\unsure{\varphi|\varnothing}\arrow[l]\arrow[d]\\S\&X_2\arrow[l]
\end{tikzcd}\right),\)}
la loi fonctorielle étant \unsure{évidente|ici évidente}.

Ce n'est pas encore la fin du \unsure{sorite}\,! Il faut \inserted{encore} interpréter le foncteur composé
\[
\mathrm{Ho}_W(\Psi)\mathrel{\overset{i_*}{\unsure{-|\to}}}\mathrm{Ho}_W(\Psi^{\sim})\xrightarrow{\deleted{\unsure{i}}\varepsilon_3^*}\mathrm{Ho}_W(e)=\mathrm{Ho}_W,
\]
où
\diagram{\begin{tikzcd}[ampersand replacement=\&,row sep=small,column sep=small]
\&\&1\arrow[dl,leftrightarrow]\&\\
{\varepsilon_3\colon \unsure{e}\simeq\{3\}}\arrow[r]\&3\&\&0\arrow[ul]\arrow[dl]\\
\&\&2\arrow[ul]\&
\end{tikzcd}}
est l'inclusion du sommet 3 dans $\Psi^{\sim}$.

Si $\mathrm{Ho}_W(\Psi)$ est interprété, conformément au lemme 1, comme $\mathscr F W_{\mathscr F}^{-1}$, ce foncteur composé n'est autre que
\diagram{\begin{tikzcd}[ampersand replacement=\&,row sep=small,column sep=small]
X_1\arrow[d,dash]\&\&{\longmapsto}\&X_1\times_S X_2\\
S\&X_2\arrow[l]\&\&
\end{tikzcd}}
\diagram{\(\displaystyle\biggl(\text{pour }\left(\begin{tikzcd}[ampersand replacement=\&,row sep=small,column sep=small]
X_1\arrow[d,"{p_1}"']\&\\S\&X_2\arrow[l,"{p_2}"]
\end{tikzcd}\right)\in\mathrm{Ob}\,\mathscr F,\quad\text{i.e. }p_1,p_2\in\mathrm{Fib}_W\biggr).\)}

On veut interpréter ce foncteur comme une $\varprojlim$ $W$-homotopique, \deleted{plus \unsure{précisé|précisément}} \inserted{i.e. \inserted{ici} comme} un \deleted{\unsure{produit|foncteur}} ``produit fibré $W$-homotopique''. De façon précise, \deleted{\unsure{on}} \unsure{c'est-|ceci|ce qui}\unsure{=|-}
% archive sheet number: 305
\scanpage{306}
\mspage{256}

\uline{Cor. 1} Considérons $p\colon\Psi\mathrel{\unsure{-|\to}}e$, et le foncteur
\[
p_\Psi^*\colon\mathrm{Ho}_W(e)=\mathrm{Ho}_W\longrightarrow\mathrm{Ho}_W(\Psi),
\]
déduit par passage aux cat. de fractions du foncteur
\diagram{\(\displaystyle\begin{aligned}
M&\longrightarrow M(\Psi)\\
X&\longmapsto\left(\begin{tikzcd}[ampersand replacement=\&,row sep=small,column sep=small]
X\arrow[d]\&\\\deleted{S}X\&X\arrow[l]
\end{tikzcd}\right).
\end{aligned}\)}

Ce foncteur admet \deleted{\unsure{du}} le \deleted{adjoint \unsure{à}} foncteur $\varepsilon_3^*\unsure{i|e}_*$ précédent comme adjoint à droite, i.e.
\[
p_{\Psi*}=\varepsilon_3^*i_*\colon\mathrm{Ho}_W(\Psi)\longrightarrow\mathrm{Ho}_W.
\]

\uline{Dém} On peut regarder $p_\Psi$ comme le composé
\[
\Psi\mathrel{\overset{i}{\unsure{-|\longrightarrow}}}\Psi^{\sim}\xrightarrow{p_{\widetilde\Psi}}e,
\]
et on a donc
\[
p_\Psi^*\simeq i^*(p_{\widetilde\Psi})^*.
\]

Donc l'existence et la détermination d'un adjoint à droite pour $p_\Psi^*$ \unsure{résultent de|résulteront de} \deleted{\unsure{de}}l'existence et de la détermination d'adj. à droite $i_*,(p_{\widetilde\Psi})_*$ de $i^*,(p_{\widetilde\Psi})^*$, via
\[
(p_\Psi)_*\simeq(p_{\widetilde\Psi})_*i_*.
\]
% archive sheet number: 306
\scanpage{307}
\mspage{257}

Donc le cor.~1 résulte \unsure{de} la formule
\[
p_{\widetilde\Psi*}=\varepsilon_3^*,
\]
qui \deleted{résulte} elle-même est un cas particulier du

\uline{Lemme} Soit $I$ dans $\mathrm{Cat}$, avec objet final $e$, considérons l'inclusion
\[
\varepsilon\colon e\longrightarrow I,\quad\text{et }p\colon I\mathrel{\unsure{-|\to}}e
\]

Soit $(M,W)$ une catégorie de modèles, et considérons
\[
\deleted{\unsure{M}}\ \mathrm{Ho}_W(I)\xleftarrow{p^*}\mathrm{Ho}_W .
\]
Ce foncteur admet un adjoint à gauche $p_*$, lequel est isom. à $\varepsilon^*$.

\uline{Dém} Considérons le couple de foncteurs adjoints
\diagram{\begin{tikzcd}[ampersand replacement=\&]
M(I)\arrow[r,shift right,"{p_*=\varprojlim_I}"']\&M\arrow[l,shift right,"{p^*}"']
\end{tikzcd}\ ;}
comme $I$ admet \deleted{\unsure{de}} l'objet final $e$, le foncteur $p_*$ est isomorphe à $\varepsilon^*$. Ces foncteurs $p^*,\varepsilon^*$ sont compatibles aux localiseurs\unsure{,|} $W(I),W$, donc on conclut par le lemme page 248.

Il ne reste plus qu'à prouver les lemmes 1 et 2. Je le ferai demain.
% archive sheet number: 307
\scanpage{308}
\mspage{258}

Aujourd'hui, j'ai \unsure{envie encore|aussi envie|très envie} de donner le \unsure{sorite} des carrés $W$-cartésiens, dans une

\uline{Proposition 19}\,: Soit
\diagram{\begin{tikzcd}[ampersand replacement=\&,row sep=small,column sep=small]
X_1\arrow[d,"{p_1}"']\&X_3\arrow[l,"{f_1}"']\arrow[d,"{f_2}"]\\
S\&X_2\arrow[l,"{p_2}"]
\end{tikzcd}}
un carré comm\unsure{:|.|. $Q$} dans $M$. Considérons des factorisations de $p_1,p_2$ en
\[
X_1\mathrel{\overset{i_1}{\unsure{-|\to}}}\overline X_1\mathrel{\overset{\bar p_1}{\unsure{-|\to}}}S
\qquad X_2\mathrel{\overset{i_2}{\unsure{-|\to}}}\overline X_2\mathrel{\overset{\bar p_2}{\unsure{\leftarrow|-|\to}}}S
\quad\left|\begin{aligned}i_1,i_2&\in W\\ \bar p_1,\bar p_2&\in F_W\end{aligned}\right.
\]
et considérons les flèches induites par $X_3\to X_1\times_S X_2$ (\unsure{déduite} de $Q$) via
\diagram{\begin{tikzcd}[ampersand replacement=\&,row sep=small]
\&\overline X_1\times_S X_2\\
X_1\times_S X_2\arrow[ur]\arrow[r]\arrow[dr]\&\overline X_1\times_S\overline X_2\\
\&X_1\times_S\overline X_2
\end{tikzcd},}
de sorte qu'on trouve \deleted{trois flèches} un diagramme comm\unsure{.}
\diagram{\begin{tikzcd}[ampersand replacement=\&,row sep=small]
\&\overline X_1\times_S X_2\arrow[d,"{j_2}"']\&\deleted{\unsure{\xi}}\\
X_3\arrow[ur,"{\pi_1}"]\arrow[r,"{\pi_0}"]\arrow[dr,"{\pi_2}"']\&\overline X_1\times_S\overline X_2\&\\
\&X_1\times\overline X_2\arrow[u,"{j_1}"]\&
\end{tikzcd},}
où $j_1=i_1\times_S\mathrm{id}_{\overline X_2}$ et $j_2=\mathrm{id}_{\overline X_1}\mathbin{\times}i_2$ sont dans $W$ (car $i_1$ et $i_2$ le sont, et $\overline X_1$ et $\overline X_2$ sont dans $\underline{\mathrm{Fib}}_W S$).

On a donc
\[
(\pi_1\in W)\Longleftrightarrow(\pi_2\in W)\qquad(\unsure{\Longleftrightarrow|\Longrightarrow}\pi_0\in W)
\]
ce qui prouve que la condition $\pi_1\in W$, qui semble dépendre du choix de la factorisation par $i_1$, n'en \uline{dépend pas}, et idem pour $\pi_2\in W$, et \deleted{\unsure{p}}pour $\unsure{\pi|\pi_0}\in W$ (qui semble dépendre à la fois de $i_1$, et de $i_2$).
% archive sheet number: 308
\scanpage{309}
\mspage{259}

\uline{Définition} Quand ces conditions sont satisfaites, on dit que le carré $Q$ est $W$-cartésien.

\uline{Corollaire 1} \deleted{Supposons que $W$ soit \unsure{forte-}} \inserted{Si $Q$ est $W$-cartésien, alors} \inserted{l'image de} $Q$ \inserted{dans $\mathrm{Ho}_W(\Psi^{\sim})$} est dans l'image essentielle de \deleted{\unsure{$\varepsilon$|$i$}}
\[
\unsure{j|i}_*\colon\mathrm{Ho}_W(\Psi)\longrightarrow\mathrm{Ho}_W(\Psi^{\sim}),
\]
i.e. la flèche canonique dans $\mathrm{Ho}_W(\Psi^{\sim})$
\[
\deleted{\unsure{i_*|i^*}}\,Q\longrightarrow i_*i^*Q
\]
est iso. La réciproque est vraie si $W$ est fortement \unsure{saturée|saturé}.

\uline{Corollaire 2} Considérons un diagramme comm. dans $M$
\diagram{\(\text{(D)}\quad\begin{tikzcd}[ampersand replacement=\&,row sep=small,column sep=small]
X'_0\arrow[dd]\&\&X'_1\arrow[ll]\arrow[dd]\&\&X'_2\arrow[ll]\arrow[dd]\\
\&1\&\&2\&\\
X_0\&\&X_1\arrow[ll]\&\&X_2\arrow[ll]
\end{tikzcd},\)}
d'où deux carrés commutatifs 1, 2. et un carré composé \unsure{3},
\diagram{\begin{tikzcd}[ampersand replacement=\&,row sep=small,column sep=small]
X'_0\arrow[dd]\&\&X'_2\arrow[ll]\arrow[dd]\\
\&3\&\\
X_0\&\&X_2\arrow[ll]
\end{tikzcd}.}

Supposons le carré 1 $W$-cartésien. Alors le carré 2 est $W$-cartésien ssi le carré 3 l'est.

\uline{Dém} Considérons une factorisation de
\diagram{\begin{tikzcd}[ampersand replacement=\&,row sep=small]
X'_0\arrow[d]\\X_0
\end{tikzcd}}
en
\diagram{\begin{tikzcd}[ampersand replacement=\&,row sep=small,column sep=small]
\overline X'_0\arrow[dr]\&X'_0\arrow[l]\arrow[d]\\
\&X_0
\end{tikzcd}}
avec $(X'_0\mathrel{\unsure{-|\to}}\overline X'_0)\in W$, et $(X'_0\mathrel{\unsure{-|\to}}X_0)\in\mathrm{Fib}_W$. On peut \unsure{donc|alors} remplacer dans (D) $X'_0$ par $\overline X'_0$, \deleted{\unsure{en s}} obtenant un diagramme $\overline D$, et on voit \unsure{aussitôt}
% archive sheet number: 309
\scanpage{310}
\mspage{260}
que \inserted{chacun} des carrés $1(D),2(D),3(D)$ \deleted{\unsure{sont}}est $W$-cartésien\unsure{s} ssi les carrés correspondants $1(\overline D),2(\overline D)$\unsure{, --|,} $3(\overline D)$ l'est. Cela montre qu'$\unsure{\mathrm{OPS}|\mathrm{SPS}}$ $X_0\mathrel{\unsure{-|\to}} S$ dans $\mathrm{Fib}_W$. Considérons alors le diagramme \inserted{comm.} déduit de $D$
\diagram{\begin{tikzcd}[ampersand replacement=\&,column sep=small]
X'_0\arrow[d]\&{\widetilde X'_{\unsure{p|1}}\overset{\mathrm{def}}{=}X'_0\times_{X_0} X_1}\arrow[l]\arrow[d]\&X'_1\arrow[l,"\alpha"']\&X'_2\arrow[l]\arrow[d]\\
X_0\&X_1\arrow[l,arrows={-{Bar}}]\&\&X_2\arrow[ll]\arrow[ll,bend right=5,dash]
\end{tikzcd},}
où $\alpha\in W$ en vertu de l'hypothèse $W$-cartésienne sur le carré 1. On voit alors que le carré 2 est $W$-cartésien ssi le carré correspondant avec $X'_1$ remplacé par $\widetilde X'_1$ l'est, \unsure{quant} aux carrés 1, 3 ils ne changent pas par cette modification. Cela montre qu'$\mathrm{OPS}$ le carré 1 \uline{cartésien}. Considérons alors le diagramme
\diagram{\begin{tikzcd}[ampersand replacement=\&,column sep=small,row sep=small]
X'_0\arrow[dd,"{\mathrm{Fib}_W}"']\&X'_1\arrow[l]\arrow[dd,"{\mathrm{Fib}_W}"]\&{\widetilde X'_2\overset{\mathrm{def}}{=}X'_1\times_{X_{\unsure{1|0}}}X_2}\arrow[l,dash]\arrow[dd,"{\mathrm{Fib}_W}"]\&X'_2\arrow[l,"\beta"']\\
\arrow[r,phantom,"{\mathrm{cart}}" description]\&\arrow[r,phantom,"{\mathrm{cart}}" description]\&\&\\
X_0\&X_1\arrow[l]\&X_2\arrow[l]\&
\end{tikzcd}.}

Alors la prop. 19 et la \deleted{\unsure{dé}}définition montrent que le carré 2 est $W$-cartésien ssi $\beta\in W$, et \unsure{de même|donc} pour le carré 3. Ce qui achève la démonstration.

\uline{Corollaire 2} Soit
\diagram{\begin{tikzcd}[ampersand replacement=\&,row sep=small,column sep=small]
X_{\unsure{1|4}}\arrow[d,"{p_1}"']\&X_3\arrow[l,"{f_2}"']\arrow[d,"{f_1}"]\\
S\&X_2\arrow[l,"{p_2}"]
\end{tikzcd}}
un carré $W$-cartésien. \deleted{Si $p_1\in W$,} Alors
\[
\begin{aligned}
p_1\in W&\Longrightarrow f_1\in W\\
p_2\in W&\Longrightarrow f_2\in W
\end{aligned}
\]

\margin{plus précisément, si $p_{\unsure{i|j}}\in W$, alors $Q$ est $W$-cartésien ssi $f_{\unsure{i|i'}}\in W$}
% archive sheet number: 310
\scanpage{311}
\mspage{261}

Supposons $p_1\in W$. Alors dans la factorisation de $p_1$ en $X_1\mathrel{-}\overline X_{\unsure{1|2}}\mathrel{\unsure{-|\to}} S$, on peut prendre $\overline X_1=S$, donc le carré est $W$-cartésien ssi
\[
(X_3\longrightarrow\overline{\deleted{\unsure{S|X}}X}_1\times_S X_2=X_2)\ \in W
\]
i.e. $f_1\in W$, qed.

Bien sûr, on ne peut en général affirmer les implications inverses ; $f_i\in W$ n'implique nullement $p_i\in W$.

Tous ces résultats, et la notion de carré $W$-cartésien, se dualisent (moyennant l'axiome \inserted{idoine} de factorisation à droite, cf page 240) en une théorie des sommes amalgamées $W$-homotopiques, et des carrés $W$-cocartésiens. Cette fois on travaille sur la catégorie
\diagram{\(\displaystyle
\unsure{\Phi|\Psi}=\left(\begin{tikzcd}[ampersand replacement=\&,column sep=tiny,row sep=tiny]
\&1\arrow[dl]\\
0\&\\
\&2\arrow[ul]
\end{tikzcd}\right)
\overset{j}{\lhook\joinrel\longrightarrow}
\Phi^{\sim}=\left(\begin{tikzcd}[ampersand replacement=\&,column sep=tiny,row sep=tiny]
\&1\arrow[dl]\&\\
0\&\&3\arrow[ul]\arrow[dl]\\
\&2\arrow[ul,dash]\&
\end{tikzcd}\right),\)}
$j$ étant une immersion fermée. \unsure{On a|Ceci|Cela} un foncteur restriction
\[
j^*:\mathrm{Ho}_W(\Phi^{\sim})\mathrel{\unsure{-|\to}}\mathrm{Ho}_W(\Phi),
\]
dont on construit un adjoint à \unsure{gauche|droite}
\[
j_{\unsure{!|*}}:\mathrm{Ho}_W(\Phi)\mathrel{\unsure{-|\to}}\mathrm{Ho}_W(\Phi^{\sim}).
\]
on trouve \deleted{\unsure{une}} notion \unsure{naturelle} de
% archive sheet number: 311
